\subsection{Finite improvement of the box exponent}
\label{sec:area_law_exponent_bootstrap}

\begin{definition}[Fixed parameters for exponent improvement]
    \label{def:area_law_bootstrap_parameters}
    \lean{
        TNLean.PEPS.AreaLaw.Scan.BootstrapParameters.g0,
        TNLean.PEPS.AreaLaw.Scan.BootstrapParameters.mu,
        TNLean.PEPS.AreaLaw.Scan.BootstrapParameters.nu,
        TNLean.PEPS.AreaLaw.Scan.BootstrapParameters.ell,
        TNLean.PEPS.AreaLaw.Scan.BootstrapParameters.kappa
    }
    \leanok
    \uses{thm:area_law_exponent_gaps}
    For $e_0\in\mathbb R$, set $e_*=2\cdot10^{-6}$; the parameters used in
    the proof of \cite[Proposition~9.5]{OpenAI2026AreaLaw} are
    \begin{align}
        g_0  & =\frac{1-e_0}{2}, & \mu    & =1-g_0,              & \nu & =\frac{g_0}{1000},
        \label{eq:area_law_bootstrap_parameters}                                            \\
        \ell & =\frac{\nu}{100}, & \kappa & =\min\{\ell,e_*/4\}.
        \label{eq:area_law_bootstrap_separation}
    \end{align}
    These parameters remain fixed as the entropy exponent decreases.
\end{definition}

\begin{theorem}[Margins in the exponent estimates]
    \label{thm:area_law_bootstrap_margins}
    \lean{
        TNLean.PEPS.AreaLaw.Scan.BootstrapParameters.parameter_bounds,
        TNLean.PEPS.AreaLaw.Scan.BootstrapParameters.density_exponents,
        TNLean.PEPS.AreaLaw.Scan.BootstrapParameters.energy_exponents,
        TNLean.PEPS.AreaLaw.Scan.BootstrapParameters.comparison_exponents
    }
    \leanok
    \uses{def:area_law_bootstrap_parameters}
    With the parameters of (\ref{eq:area_law_bootstrap_parameters}) and
    (\ref{eq:area_law_bootstrap_separation}), the condition $0<e_0<1$ gives
    $0<g_0<1/2$, $0<\ell<1$, and $\kappa>0$.
    If $e_0<1$ and $e\le e_0$, the density-error exponents satisfy
    \begin{align}
        e_0-\mu+\nu & =-0.999g_0,
                    & -\frac{g_0-\ell}{4}+\nu & =-0.2489975g_0,
        \label{eq:area_law_bootstrap_density_values}            \\
        e-\mu+\nu   & <-0.24g_0,
                    & -\frac{g_0-\ell}{4}+\nu & <-0.24g_0.
        \label{eq:area_law_bootstrap_density_margin}
    \end{align}
    For $e_0<1$, write $p=\ell+2\kappa$ for the energy-prefactor exponent. Then
    \begin{align}
        p & \le0.00003g_0, & p-\nu & <-\nu/2, & p-0.03g_0 & <-\nu/2.
        \label{eq:area_law_bootstrap_energy_margin}
    \end{align}
    For $e_0>0$ and $e>e_*$, one has $\kappa<(1-\ell)e$ and
    $g_0-\ell-\nu/4<1$.
\end{theorem}
\begin{proof}\leanok
    If $0<e_0<1$, then $g_0=(1-e_0)/2$ lies in $(0,1/2)$. Hence
    $\ell=g_0/100000$ lies in $(0,1)$, and $\kappa=\min\{\ell,e_*/4\}>0$.
    Substituting (\ref{eq:area_law_bootstrap_parameters}) and
    (\ref{eq:area_law_bootstrap_separation}) gives
    (\ref{eq:area_law_bootstrap_density_values}). Since $g_0>0$ when $e_0<1$,
    the identities imply (\ref{eq:area_law_bootstrap_density_margin}) for
    $e\le e_0$. The bound $\kappa\le\ell$ gives
    (\ref{eq:area_law_bootstrap_energy_margin}). For the comparison bounds,
    substitute the same formulas and use $\kappa\le e_*/4$, $e>e_*$, and
    $e_0>0$.
\end{proof}

\begin{theorem}[Binary-entropy cost of the rounded bands]
    \label{thm:area_law_bootstrap_binary_entropy_cost}
    \lean{TNLean.PEPS.AreaLaw.Scan.exists_bootstrap_binaryEntropy_cost_bound}
    \leanok
    Fix the scanner exponents described in Section~\ref{sec:al_scanner}, and put
    \begin{align}
        L_n & =\lfloor n^{1-\ell}\rfloor,\qquad
        m_n=\lfloor n^\mu\rfloor,\qquad
        K_n=\left\lfloor\frac{L_n}{8m_n}\right\rfloor,\qquad
        a_n=\frac4{K_n}.
    \end{align}
    There is an integer $N\ge1$, depending only on these fixed exponents,
    such that $K_n\ge1$ and $a_n>0$ for every integer $n\ge N$.
    Define $h(\tau)=-\tau\log|\tau|-(1-\tau)\log|1-\tau|$ for
    $\tau\in\mathbb R$, with the convention $0\log0=0$.
    On $[0,1]$, this is the binary entropy.
    For every such $n$ and every $\tau\in\mathbb R$,
    \begin{align}
        \frac{h(\tau)}{a_n}\le\frac{\log2}{4}\,n.
    \end{align}
    The coefficient is absolute, and the threshold precedes $\tau$ and
    every physical instance. This bounds the binary-entropy term in the
    proof of \cite[Proposition~9.5]{OpenAI2026AreaLaw} for the actual
    rounded band count.
\end{theorem}
\begin{proof}\leanok
    Since $\mu<1-\ell$, the numerical scale estimates give $K_n\ge1$
    for all sufficiently large $n$. They also give $L_n\le n$, and
    integer division gives $K_n\le L_n$. The extension satisfies
    $h(\tau)\le\log2$ for every real $\tau$, so
    \begin{align}
        \frac{h(\tau)}{a_n}
        =\frac{K_nh(\tau)}4
        \le\frac{K_n\log2}4
        \le\frac{n\log2}4.
    \end{align}
\end{proof}

% The comparison-budget statement and proof are still under construction.
\begin{theorem}[Additive comparison costs at fixed bootstrap scales]
    \label{thm:area_law_bootstrap_comparison_budget}
    \lean{TNLean.PEPS.AreaLaw.Scan.exists_bootstrap_comparison_budget_bound}
    \uses{def:area_law_bootstrap_parameters}
    Fix $e_*<e_0<1$, where $e_*=2\cdot10^{-6}$, and use the parameters of
    (\ref{eq:area_law_bootstrap_parameters}) and
    (\ref{eq:area_law_bootstrap_separation}). Set $\eta=1-\ell$ and
    \begin{align}
        L_n & =\lfloor n^\eta\rfloor,\qquad
        m_n=\lfloor n^\mu\rfloor,\qquad
        K_n=\left\lfloor\frac{L_n}{8m_n}\right\rfloor,\qquad
        D_n=\lceil n^\kappa\rceil,\qquad a_n=\frac4{K_n}.
    \end{align}
    There is an integer $N\ge1$, depending only on $e_0$, such that
    $K_n\ge1$ and $a_n>0$ for every integer $n\ge N$ and, simultaneously,
    \begin{align}
        nL_n^e+nD_n+n^{3/5}+\frac{h(\tau)}{a_n}+1
        \le\left(5+\frac{\log2}{4}\right)n^{1+\eta e}
    \end{align}
    for every $e_*<e\le e_0$ and every real $\tau$, where
    $h(\tau)=-\tau\log|\tau|-(1-\tau)\log|1-\tau|$ with the convention
    $0\log0=0$. The threshold is chosen before $e$ and $\tau$.
    This is a numerical estimate for the additive terms following the
    comparison in the proof of \cite[Proposition~9.5]{OpenAI2026AreaLaw}.
\end{theorem}

\begin{theorem}[Uniform bounds after finitely many exponent improvements]
    \label{thm:area_law_finite_exponent_improvement}
    \lean{TNLean.PEPS.AreaLaw.Scan.exists_uniform_boxError_bound_of_eventual_improvement}
    \leanok
    \uses{def:area_law_bootstrap_parameters}
    Fix $0<e_0<1$, set $e_*=2\cdot10^{-6}$, and take $\ell$ from
    (\ref{eq:area_law_bootstrap_separation}).
    Let $F_i:\mathbb N\to\mathbb R$ be an arbitrary family, and suppose there
    are constants $V\ge0$ and $C_0>0$ such that, for every $i$ and $r\ge1$,
    \begin{align}
        F_i(r) & \le Vr^2, & F_i(r) & \le C_0r^{1+e_0}.
        \label{eq:area_law_bootstrap_initial}
    \end{align}
    Assume the following implication for every $e_*<e\le e_0$: if there is a
    constant $C_e>0$ such that $F_i(r)\le C_er^{1+e}$ for every $i$ and
    $r\ge1$, then there are a constant $C'_e>0$ and an integer $N_e$ such that
    \begin{align}
        F_i(r) & \le C'_er^{1+(1-\ell)e}
        \label{eq:area_law_bootstrap_step}
    \end{align}
    for every $i$ and $r\ge N_e$. Then there is one constant $C>0$ such that
    $F_i(r)\le Cr^{1+e_*}$ for every $i$ and every $r\ge1$.
\end{theorem}
\begin{proof}\leanok
    \uses{thm:area_law_bootstrap_margins}
    Write $a=1-\ell$, so $0<a<1$. At each improvement, enlarge the constant
    using the quadratic bound in (\ref{eq:area_law_bootstrap_initial}) to include
    $1\le r<N_e$. The enlarged constant is independent of $i$.
    Starting at $e_0$, apply (\ref{eq:area_law_bootstrap_step}) while the current
    exponent exceeds $e_*$. Since $a^ne_0\to0$, only finitely many steps are
    needed. Once $a^ne_0\le e_*$, the inequality
    $r^{1+a^ne_0}\le r^{1+e_*}$ for $r\ge1$ gives the result.
\end{proof}

% This is the conditional finite iteration in the proof of Proposition 9.5.
% Its one-step analytic implication remains a hypothesis; the source
% proposition is not marked as proved by this entry.
